\section{Discussion}

We now discuss the results and present a model for the photoionization and recombination process. First, we develop a model for the photoionization and recombination process. Then, we present an assessment of the viability of the photoionization process to produce a net return in energy in seeded oxy-fuel combustion products.

\subsection{PhotoIonization Process}

We propose that photoionization occurs as the following two step process. 

\begin{enumerate}
\item KOH + $h\nu$ \rightarrow K* + OH
\item K* + M \rightarrow K+ + e- + M
\end{enumerate}



Step 1 is the first step of the chemiluminescence process discussed in the methods section. The photodiode signal is also is a similar measure of the KOH chemiluminescence and is observed to scale with the nominal potassium concentration as presented in Figure\ref{fig:kwt_ionization}B. The $\Delta AS_{max}$ also shown in Figure\ref{fig:kwt_ionization}B is observed to scale with the nominal potassium concentration. Assuming that $\Delta AS_{max}$ is related to the laser-induced free electron population, this observation is evidence for a laser-induced ionization process that occurs in addition to the laser-induced chemiluminescence proposed in step 2. 

It is possible that K atoms are directly ionized in this process. However, the UV absorption cross section of KOH is much larger than other major species present including K (see SI). Additionally, the AS signal linearly increases with the nominal potassium concentration, even though the atomic potassium density measurements at the goldilocks region decrease with nominal potassium mass fraction (\ref{fig:kwt_ionization}A). However, the predicted KOH concentration at the goldilocks zone from CFD scales linearly with the nominal Kwt along with the PD and $\Delta AS_{max}$ signals. As mentioned previously and shown in Figure \ref{fig:kwt_ionization}A, this KOH concentration is nearly equal to the measured K concentration at the barrel exit, implying that most atomic K has become bound in KOH by the time it reaches the goldilocks region.

Together, this evidence suggests an ionization process that is related to the KOH chemiluminescence process, and not the direct photoionization of atomic potassium. In this step the excited K atom might be ionized by a collision with a third body. According to the ladder model developed by Alkemade et al. (discussed in the introduction), the K +M \rightarrow K+ + e- +M process depends on a collision with a third body, and appears to be dependent on the potassium atom being in an excited state.[Alkemade] It is also possible that the excited K atom is ionized by additional UV photons. The UV cross section of K is much smaller than KOH, but it is possible that the 2P0 state of K has a much higher UV cross section (see SI). The quadratic power dependence of the $\Delta AS_{max}$ signal (Figure \ref*{fig:SI_power_dependence}) may indicate a two-body process. 

The $\Delta AS_{max}$ signal also is observed to peak downstream of the barrel exit. A small peak in a similar location is observed in the CFD KOH concentration, but generally the centerline CFD-predicted KOH concentration is weakly dependent on position and not sufficient to explain the peak in the $\Delta AS_{max}$ signal (Figure \ref*{fig:SI_5x3_pos_mws_KOH}). We do observe from the ICCD camera images for different positions that the extent of the chemiluminescence is smaller near the barrel exit and mirrors the KOH spreading. We suspect that the change in microwave transmission is related the physical size of the laser-jet interaction and could explain the peak in $\Delta AS$ signal. In this work, we do not model the transmission of microwaves between the two horns, and further work is planned to account for this effect.

\subsection{Recombination Process}

After ionization, the excited electron population decays back to equilibrium through a recombination process. For the low temperature conditions of the SFR-maximized position, we propose that molecular oxygen is likely the dominant recombination partner

\begin{enumerate}
\item e- + O2 + M \rightarrow O2- + M 
\end{enumerate}

The measured exponential time constant of the AS signal is within the range of the literature values for the predicted exponential time constant of this process as shown in Figure \ref{fig:kwt_recombination}. Axford 1997, the latest source on this process, concluded with a value ($O2_A$ in Figure \ref{fig:kwt_recombination}) of $5 \times 10^{-34} cm^3/s$, which gives values about an order of magnitude lower than our experimental measurements. A negative temperature coefficient for $k_{r,b,O2}$ as proposed by Goodings could explain the discrepancy. However, the actual $k_{r,b,O2}$ measured by Goodings ($O_{2,C}$) is significantly longer than the measured exponential time constant. 

Another piece of evidence for this recombination process is that the exponential time constant is observed to be somewhat independent of the position along the free jet (Figure \ref*{fig:SI_516_mws_decaytime}). The exponential time constant is difficult to quantify outside of the SFR-maximized region due to the AS fluctuations, but the AS magnitude is still appreciable after a few microseconds regardless of position. This is a similar order of magnitude to the exponential fit time constants obtained at the SFR-maximized position. CFD simulations (SI TODO) show that the oxygen concentration is relatively constant along the free jet. This constant oxygen concentration occurs from running fuel lean, and air entrainment.

Finally, we have observed dependence of exponential time constant on equivalence ratio (Figure \ref*{fig:SI_516_mws_decaytime}). However, similar to the position dependence, quantification of the exponential time constant is difficult due to the AS fluctuations. Additionally, the profile of the torch changes with equivalence ratio (extends for high equivalence ratio), so it is difficult to deconvolve the effects of the profile change from the effects of the equivalence ratio.



\subsection{Viability Analysis}

With the developed understanding of the photoionization and recombination process, we performed an assessment of the viability of the photoionization process to produce a net return in energy in seeded oxy-fuel combustion products. The methods are described in detail in the SI and outlined here. We performed with 0D chemical equilibrium calculations to determine the composition and thermophysical properties of a seeded oxy-kerosene combustion products at various temperatures and pressures, approximating various locations in a MHD generator. We then apply a nonequilibrium ionization model to determine the change in electrical conductivity due to photoionization. We then calculate the MHD power output for various photoionization power inputs and determine the regimes in which photoionization is a net energy return. We calculate a figure of merit for the photoionization process, $\gamma$, which is the ratio of the change in MHD power output, $dP_{MHD}$ for an intestinal amount of photoionization power input, $dP_{in}$. Therefore, a $\gamma > 1$ indicates a net energy return. We note that a $\gamma > 1$ does not necessarily indicate the viability of the overall MHD power generation process, just that the photoionization process is a net energy return. We derive an expression for $\gamma$ in the SI, and it is given by


\begin{equation}
\gamma = \frac{dP_{MHD}}{dP_{in}} \Big|_{P_{in}=0} = \frac{e \mu_e u^2 B^2 \eta}{E_{IP}k_{r, m, eff}} 
\end{equation}


For the MHD power output, we use the ideal power density in a MHD generator, given by \ref{eq:mhd_ideal_power}. $E_{IP}$ is the ionization potential of the seed species, $k_{r, m, eff}$ is the effective recombination rate, and $\eta$ is the photoionization efficiency. $\eta$ describes the fraction of the photoionization power input that is converted to ionized electrons, which each require $E_{IP}$. We analyze two cases for $\eta$, a perfect efficiency case, $\eta=1$, and a case where the photoionization efficiency is determined by the fraction of UV light absorbed by the KOH, $\eta_{KOH}$. 

Figure \ref{fig:viability_gamma}A shows $\gamma$ for various temperatures and pressures for the case of $\eta=1$. The assumed values for other parameters are indicated in the caption.  For $\gamma > 1$ at low temperatures and pressures (blue region), the photoionization process exhibits a net energy return and is viable. For $\gamma < 1$, at high temperatures and pressures (red region) the photoionization process is not viable. The line $\gamma = 1$ is the boundary between these two regions and is defined as the 'viability line'. To determine the viability line (green dashed line), we interpolate $\gamma (P)$ each temperature to determine the pressure at which $\gamma = 1$. 

In Figure \ref{fig:viability_gamma}B, we show the viability lines for some different analysis conditions. The effect of equivalence ratio is shown, where fuel lean mixtures are less viable at low temperatures, but the equivalence ratio has little effect at high temperatures. This is due to a change in dominant recombination partner. In Figure S\ref{fig:SI_P_zero_rxn_component} we examine viability lines considering only one reaction partner (e.g. only O2, only K+, etc.) to determine the effect of each reaction on the shape of the viability line. We observe that for temperatures below about 2000K, O2 is the dominant recombination partner for fuel lean mixtures and K+ is the dominant recombination partner for fuel rich mixtures. For higher temperatures, the temperature dependence of the H2O reaction causes it to dominate the recombination process, though recombination by K+ is still significant. Figure S\ref{fig:SI_krm_cfd_pos} shows the predicted lifetimes from different species for different axial position in the free jet and predicts that H2O and K+ overtake O2 as the dominant recombination partner at the higher temperatures closer to the barrel exit. However, as mentioned previously, the MWS SFR closer to the barrel exit is too high to accurately measure the exponential time constant and validate these predictions. Further work is needed to develop a better understanding of the recombination process at these high temperatures.  

The difference between the perfect efficiency case and the KOH absorption efficiency case is also shown in Figure \ref{fig:viability_gamma}B. As noted previously and shown in Figure \ref{fig:SI_UV_abscs} the UV absorption cross section of KOH is much larger than other major species present. In Figure \ref*{fig:SI_atten_length} we calculate the fraction of light absorbed by KOH and find that for temperatures below about 2500K, nearly all the UV light is absorbed by KOH, with a decrease to about 10\% around 3700 K. Therefore, there is not a significant difference between the perfect efficiency case and the KOH absorption efficiency case for temperatures below 2500K. A perhaps more important quantity would be the yield of excited KOH to ionized electrons. We have explored fitting of the non-exponential part of the MWS decay to determine the yield of excited KOH to ionized electrons, but further work is needed to understand the accuracy of these results. In any case, the $\eta_{perf}$ and $\eta_{KOH}$ cases are upper bounds on the photoionization viability.  

We also analyze the targeting of electrode boundary layers in our viability model. As discussed in the introduction, if the photoionization power input only needs to be applied to a small volume of the generator to selectively target cold boundary layers near the electrode surfaces, the photoionization process can be more viable than attempting to enhance the conductivity of the entire channel. We consider MHD channel comprising a bulk region and a boundary layer layer region in series, and the bulk region width being a fraction $l_b$ of the overall channel width. The bulk region is fixed at 3000 K and the temperature of the boundary layer region is varied. The photoionization power input is applied only to the boundary layer region, and the change in overall channel conductivity is calculated with a 1D series resistance model. Figure \ref{fig:viability_gamma}B shows the viability lines for two different $l_b$ values. For low temperatures, the photoionization process becomes significantly more viable for low $l_b$ values, with the viability at atmospheric pressure becoming well within the viability range for approximately 2000 K and below. Near and above the fixed bulk temperature of 3000K, the viability actually decreases for boundary layer targeting. This is due to the fact that energy inputs to the boundary layer are not as effective at increasing the overall channel conductivity. However, the channel conductivity and MHD generator performance is already high at these temperatures. 


\begin{figure}[h]
    \centering
    \includegraphics[width=\linewidth]{\repodir/final/figures/output/Fig7_Viability.png} 
    \caption{Figure of merit for photoionization for various temperatures and pressures. The lines represent where $\gamma = 1$ for various temperatures and pressures as described in the text.   The plots are all for 1\% K mass fraction and assuming a $u=1e3 m/s$, $B=5T$. $\mu_e$ is calculated from collision cross sections and species concentrations as described in the SI. $k_{r,m,eff}$ is calucalted as a sum over experimental O2, K+, H2O,
and OH monomolecular recombination rates. A) $\gamma$ for various temperatures and pressures for the case of $\eta=1$. B) Viability lines for different analysis conditions. The effect of equivalence ratio, $\eta$, and bulk region width $l_b$ are shown.}
    \label{fig:viability_gamma}
\end{figure}



\subsection{Conclusion}

In conclusion, we have presented experimental evidence of UV photoionization in a potassium seeded high-velocity oxy-fuel flame. The experimental results suggest and ionization mechanism of excimer-induced KOH dissociation, followed by thermal ionization of the excited K atom. For the recombination process, we have observe a first-order recombination process that is consistent with the capture of free electrons by oxygen molecules. 

Using these results, we have developed a model for the photoionization and recombination process and determined the viability of the photoionization process to produce a net return in energy in seeded oxy-fuel combustion products. The viability analysis suggests that for targeting of the entire bulk MHD channel at atmospheric pressure the enhancement from non-equilibrium excitation can be expected to be low, except perhaps at low temperatures in fuel-rich conditions. On the other hand, if the photoionization power input only needs to be applied to a small volume of the generator to, for example, selectively target cold boundary layers near the electrode surfaces, we have found that the photoionization process is significantly more viable. Overall, our analysis suggests that photoionization should be explored further as a method of selectively targeting electrode boundary layers or other problematic low-temperature regions in MHD generators.